% ============================================================
%  ANEXOS
% ============================================================

\chapter{Especificaciones técnicas de la cámara MicaSense RedEdge-MX}
\label{anx:camara}

La Tabla~\ref{tab:specs_camara} resume las especificaciones técnicas
principales de la cámara MicaSense RedEdge-MX empleada en este trabajo.

\begin{table}[H]
\centering
\caption{Especificaciones técnicas de la cámara MicaSense RedEdge-MX}
\label{tab:specs_camara}
\begin{tabular}{ll}
\toprule
\textbf{Parámetro}              & \textbf{Valor} \\
\midrule
Número de bandas                & 5 \\
Resolución por banda            & 1280 $\times$ 960 px (1.2 Mpx) \\
Tamaño del sensor               & 4.8 mm $\times$ 3.6 mm \\
Tamaño de píxel                 & 3.75 \si{\micro m} \\
Distancia focal                 & 5.5 mm \\
Campo de visión horizontal      & 47.2° \\
Campo de visión vertical        & 35.4° \\
Lentes                          & 5 objetivos físicamente separados \\
Interfaz de captura             & API HTTP REST (\texttt{192.168.10.254}) \\
Formato de salida               & TIFF sin comprimir (16 bits) \\
Sensor de irradiancia           & DLS 2 (Downwelling Light Sensor) \\
Alimentación                    & 4.2--15.8 V, consumo típico 4 W \\
Masa                            & 232 g \\
\midrule
\multicolumn{2}{c}{\textit{Bandas espectrales}} \\
\midrule
Banda 1 — Azul (Blue)           & Centro: 475 nm, FWHM: 32 nm \\
Banda 2 — Verde (Green)         & Centro: 560 nm, FWHM: 27 nm \\
Banda 3 — Rojo (Red)            & Centro: 668 nm, FWHM: 16 nm \\
Banda 4 — Red Edge              & Centro: 717 nm, FWHM: 12 nm \\
Banda 5 — NIR                   & Centro: 842 nm, FWHM: 57 nm \\
\bottomrule
\end{tabular}

  {\footnotesize Fuente: elaboración propia.}
\end{table}

% ============================================================

\chapter{Código fuente del \textit{pipeline} de procesamiento}
\label{anx:codigo}

La Tabla~\ref{tab:scripts} describe los scripts principales del
\textit{pipeline} de preprocesamiento, extracción de características y
clasificación que conforman el sistema implementado.

\begin{table}[H]
\centering
\caption{Scripts principales del \textit{pipeline} de procesamiento}
\label{tab:scripts}
\small
\begin{tabular}{p{6.0cm}p{6.5cm}}
\toprule
\textbf{Script} & \textbf{Descripción} \\
\midrule
\texttt{pipeline\_completo.py}
  & Orquestador principal. Ejecuta las tres etapas secuencialmente para una sesión de campo. \\[4pt]
\texttt{preprocesamiento.py}
  & Calibración radiométrica mediante \ac{CRP}, alineación de bandas por \ac{SIFT}+RANSAC y recorte de bordes. \\[4pt]
\texttt{extraccion\_features.py}
  & Extracción de ventanas ($128\times128$~px, 50\,\% solapamiento), cálculo de 40+ índices espectrales y 9 estadísticas por índice. \\[4pt]
\texttt{clasificador\_rf.py}
  & Entrenamiento del \ac{RF} con validación cruzada estratificada. \\[4pt]
\texttt{evaluar\_loso.py}
  & Validación \ac{LOSO} con $N$ sesiones. Genera métricas por pliegue y promedio. \\[4pt]
\texttt{extraccion\_segmentada.py}
  & Re-etiquetado de ventanas por bounding box usando anotaciones Roboflow. \\[4pt]
\texttt{metricas\_modelo\_final.py}
  & Entrenamiento y evaluación del modelo final: EasyEnsemble x10 XGBoost,
    top-70 características seleccionadas por XGBoost, validación \ac{LOSO}. \\[4pt]
\texttt{micasense\_preview\_capture.py}
  & Captura de imágenes mediante la API HTTP de la cámara MicaSense. \\
\bottomrule
\end{tabular}

  {\footnotesize Fuente: elaboración propia.}
\end{table}

% ============================================================

\chapter{Base de datos de imágenes multiespectrales}
\label{anx:dataset}

La Tabla~\ref{tab:estructura_dataset} resume la estructura del dataset
al cierre del período experimental.

\begin{table}[H]
\centering
\caption{Estructura del dataset de imágenes multiespectrales}
\label{tab:estructura_dataset}
\begin{tabular}{lc}
\toprule
\textbf{Parámetro}                      & \textbf{Valor} \\
\midrule
Sesiones oficiales (S2--S7)    & 6 \\
Zonas de mina por sesión                & 8 (2 réplicas $\times$ 4 profundidades) \\
Profundidades evaluadas                 & 1, 3, 5 y 7 cm \\
Zonas de control por sesión (S3+)       & 4 \\
Bandas espectrales por captura          & 5 (TIF, 16 bits por banda) \\
Objetos de interferencia por zona       & 3 (roca, botella de vidrio, lata) \\
Formato del dataset procesado           & CSV (una ventana por fila, 378 columnas) \\
Tamaño de ventana                       & 128 $\times$ 128 px \\
Solapamiento entre ventanas             & 50\,\% \\
\bottomrule
\end{tabular}

  {\footnotesize Fuente: elaboración propia.}
\end{table}
